% !TEX root = ../main.tex
\section{Introduction}
\label{sec:introduction}

Automation tool marketplaces---platforms where developers publish reusable
scripts, scrapers, and integrations---are being opened to a new kind of
consumer: AI agents that autonomously discover, invoke, and pay for tools, a
shift documented in census analyses of the MCP tool ecosystem
\cite{stein2026mcpecosystem} and the GPT Store \cite{su2024gptstore}. For a tool
to enter an agent's
reach, however, it must first satisfy a platform-imposed architecture. On the
Apify Store, a tool is \emph{agentic-ready} only if it runs with limited
permissions and no standby mode. Agentic readiness is thus not a label a
developer claims, but a set of architectural constraints the platform enforces
on which tools machine consumers may touch. Yet little is known about how this
readiness architecture shapes adoption and concentration across the
marketplace.

Automation tool marketplaces are a distinct platform
type with four structural features: they host \emph{single-function} tools
rather than the multi-feature applications that mobile app stores manage
through update review \cite{comino2019appstore}, use
predominantly \emph{usage-based}
pricing rather than subscription or purchase models, organize tools through
\emph{category taxonomies} for both human and programmatic discovery, and rely
on open developer-side participation that produces a long-tailed distribution
of activity. They are not app stores (tools are headless scripts, not
interactive applications), not API marketplaces (tools encapsulate orchestrated
workflows rather than exposing raw endpoints), and not integration platforms
(tools are standalone rather than connectors between services). We describe
these features concretely for Apify in \autoref{sec:context}.\label{sec:bg:automation}

We analyze the \emph{Apify Store}, a public automation marketplace hosting
67,609 tools (``actors'') built by 5,789 developers. We define a pay-per-event
(PPE) tool as \emph{agentic-ready} when it is configured with limited
permissions and no standby mode---two architectural properties that, in the
platform's design, determine whether a tool can be autonomously discovered and
invoked by agents. 50{,}917 of the 53{,}502 PPE tools with architecture data
(95\%) satisfy this definition. We study how agentic readiness is associated
with adoption and concentration, using within-developer fixed effects that
compare agentic-ready and non-agentic-ready tools published by the \emph{same}
developer.

Prior work leaves four gaps: no empirical study of automation tool marketplaces
as a distinct platform type, no quantification of the adoption gap associated
with agentic readiness, no documentation of category-level heterogeneity in
that gap, and no analysis of whether agentic readiness is associated with
marketplace concentration. Our study addresses all four.

This study's organizing question is: what adoption and concentration
differences are associated with a platform's readiness boundary---the
constraints it imposes on which tools machine consumers may use? We find three
things.

First, agentic readiness is strongly associated with adoption. Within the same
developer---specifically the 343 who publish both agentic-ready and
non-agentic-ready tools---agentic-ready tools attract a large positive mean
adoption gap ($p < 10^{-10}$): 40\% to 72\% on the
$\log(1+\text{monthly\_users})$ scale depending on weighting, and several-fold
larger on the count scale; a within-developer similarity match confirms a positive mean premium of roughly 29\%,
indicating that tool-type composition explains part of the raw gap but not all of it. Second, the association also appears on both the
extensive and conditional intensive margins: agentic-ready tools are 17
percentage points more likely to have any users at all, and, conditional on
having users, they attract 60\% more on the log-plus-one scale. Third, the
association is positive across
categories---Bonferroni-significant in 7 of 15 under primary-category assignment
and 12 of 18 under multi-hot, in tool-level regressions---but we find no robust
evidence of higher concentration: Gini coefficients are nearly identical across
agentic-ready and non-agentic-ready tools (0.935 versus 0.946), and upper-tail
differences are too imprecisely estimated to be informative.

Across these findings, the through-line is the readiness boundary: tools a
platform configures for agent-initiated use attract substantially more
adoption, without any detectable increase in concentration. This adoption
premium is large and
consistent. The observed premium attaches to platform readiness; distinguishing agent routing from correlated architectural choices requires agent-side traffic data.

To understand how agentic readiness shapes marketplace stratification, we
organize the analysis around three research questions:

\begin{description}
  \item[RQ1.] What is the adjusted association between agentic readiness and
    tool adoption, after controlling for developer heterogeneity and tool
    characteristics?
  \item[RQ2.] On which margin of adoption does the association operate---the
    extensive margin (any users) or the intensive margin (users conditional on
    traction)---and which candidate mechanisms the data can and cannot rule
    out?
  \item[RQ3.] How does the association vary across tool categories, and is
    agentic readiness associated with higher or lower within-group adoption
    inequality?
\end{description}

This study makes two contributions to HCI and platform studies:

\begin{enumerate}
  \item \textbf{Quantifying the adoption premium of agentic readiness.} Within
    the same developer, agentic-ready tools attract a large, consistently
    positive adoption gap (40\% to 72\% on the log-plus-one scale) across both
    the extensive and intensive margins, with no robust evidence of a
    corresponding increase in within-group concentration.

  \item \textbf{The admission boundary as a governance boundary.} This premium
    reveals a design-relevant empirical signal: an automation marketplace's
    permission architecture---the limited-permissions, non-standby criteria that
    gate agent access---doubles as its discovery architecture, so access rules
    allocate visibility across human users, agent users, and developers.
\end{enumerate}

We next describe the setting, data and methods, results, and their
implications for dual-audience governance.
